% ECONOMICS_PROBLEM: {"id": "C73-10", "title": "Ordinary uniform approximate equilibria in quitting games", "jel_code": "C73", "source_id": "OP-0591"}
% FIRST_POSTED: 2026-10-09
% ATTEMPT_STATUS: unresolved
% ENTRY_KIND: research_attempt
% !TEX program = pdflatex
% Self-contained source; no external BibTeX file or image assets required.
\documentclass[11pt,a4paper]{article}
\usepackage[T1]{fontenc}
\usepackage[utf8]{inputenc}
\usepackage{lmodern}
\usepackage[margin=1in]{geometry}
\usepackage{amsmath,amssymb,amsthm,mathtools}
\usepackage{booktabs,array,microtype,enumitem}
\usepackage[colorlinks=true,linkcolor=blue!55!black,citecolor=blue!55!black,urlcolor=blue!55!black]{hyperref}
\usepackage{xcolor}
\usepackage{graphicx}
\allowdisplaybreaks[2]
\setlength{\parskip}{0.3em}
\setlength{\parindent}{1.35em}
\setlist{nosep}
\numberwithin{equation}{section}
\theoremstyle{plain}
\newtheorem{theorem}{Theorem}[section]
\newtheorem{proposition}[theorem]{Proposition}
\newtheorem{lemma}[theorem]{Lemma}
\newtheorem{corollary}[theorem]{Corollary}
\theoremstyle{definition}
\newtheorem{definition}[theorem]{Definition}
\newtheorem{example}[theorem]{Example}
\theoremstyle{remark}
\newtheorem{remark}[theorem]{Remark}
\newcommand{\NN}{\mathbb N_0}
\newcommand{\EE}{\mathbb E}
\newcommand{\PP}{\mathbb P}
\newcommand{\one}{\mathbf 1}
\newcommand{\Reg}{\operatorname{Reg}}
\newcommand{\NE}{\operatorname{NE}}
\newcommand{\Delt}{\Delta}
\newcommand{\cC}{\vec C}
\newcommand{\eps}{\varepsilon}
\newcommand{\given}{\,|\,}
\newcommand{\R}{\mathbb R}
\title{Quitting Games and Uniform Equilibria:\\
\large A Consolidated Mathematical Audit of Truncation, Cyclic Equilibria, and Finite Certificates}
\author{GPT 6 Astra Pro}
\date{8 October 2026}
\begin{document}
\noindent\textbf{Problem ID:} \texttt{C73-10}\quad\textbf{Model:} GPT 6 Astra Pro\quad\textbf{Version:} 1\par
\noindent\textbf{Status:} Unresolved for the complete catalogue problem; this manuscript records partial results and research attempts.\par
\maketitle
\begin{abstract}
We investigate the open question whether every quitting game with finitely many
players admits an ordinary uniform $\eps$-equilibrium for every $\eps>0$.
We collect and carefully distinguish three kinds of conclusions. First, the
stopping-time representation yields an exact identity for the only discontinuity
created when quitting times escape to infinity, and an error bound for Nash
equilibria of finite terminal truncations. Second, in an explicit classical
three-player cyclic quitting game, we prove that the infinite game has a
subgame-perfect exact equilibrium although \emph{every} exact Nash equilibrium
of \emph{every} finite terminal truncation has unrestricted deviation gain
strictly exceeding $1/4$; consequently, a tempting universal truncation
criterion fails. Third, a quantile-rounding and date-compression argument
gives a quantitative finite-grid approximation to the infimum of unrestricted
regret, with error at most $(2n-1)D/k$. In particular, a game without
arbitrarily accurate undiscounted equilibria would admit a rigorous finite
nonexistence certificate. We supply explicit conversions to long-run-average
and discounted equilibrium inequalities and illustrate the certificates by
exact rational grid calculations. These statements do \emph{not} decide
whether every finite-player quitting game has uniform $\eps$-equilibria.
\end{abstract}

\tableofcontents

\section{The problem, literature, and scope}\label{sec:intro}
A \emph{quitting game} is a stochastic game with one nonabsorbing state.
At each stage all players simultaneously choose to continue or to quit.
The first stage at which somebody quits determines the absorbing payoff;
if everybody always continues, a specified continuation payoff applies.
Despite this spare description, the existence of ordinary approximate
equilibria for an arbitrary number of players has proved difficult.

The classical treatments of Solan and Vieille \cite{SolanVieille2001} and
Flesch, Thuijsman, and Vrieze \cite{Flesch1997} exhibit some of the
underlying phenomena. Two- and three-player quitting games admit the
relevant approximate equilibrium payoffs, while the unrestricted existence
question for four or more players is identified as open in the 2026 treatment
of Ashkenazi-Golan, Krasikov, Rainer, and Solan \cite{AKRS2026}.
Solan and Solan \cite{SolanSolan2020} prove sunspot $\eps$-equilibrium
existence for every finite-player quitting game, but public randomization
is additional structure, not an ordinary equilibrium of the game as given.
The same paper obtains a positive ordinary uniform-equilibrium conclusion
under a non-$Q$-matrix condition; other positive classes follow from
absorption-path arguments \cite{AKRS2024,AKRS2026}.

There is a distinction between exact $0$-equilibria and approximate
equilibria for every $\eps>0$. Exact $0$-equilibria need not exist
\cite{Flesch1996,AKRS2026}. Throughout, our \emph{open question} concerns
existence \emph{for each positive tolerance}, with strategies permitted to
depend on that tolerance. We do not assert that exact uniform equilibria
always exist.

Our analysis of this question developed in stages. An initial approach
tried to take exact equilibria of finite terminal games and send the
terminal date to infinity. The approach admits a clean error lemma, but
Section~\ref{sec:counterexample} disproves its universal convergence
requirement. A replacement takes \emph{arbitrary} finite-support strategies,
optimizes their regret in the \emph{infinite} game, and proves a quantitative
approximation theorem (Section~\ref{sec:finite-grid}). The final implication
is conditional and gives no positive or negative resolution of the open
question.

\section{Model, evaluations, and strategic regret}\label{sec:model}
Let $I=\{1,\ldots,n\}$ with $n\ge2$. For each subset $S\subseteq I$ let
$r^S\in\R^n$. If exactly the players in a nonempty set $S$ quit at the
first terminating stage, payoffs equal $r^S$ thereafter. Write
$c=r^\varnothing$ for the payoff vector if nobody ever quits.
We often use the convenient normalization
\begin{equation}\label{eq:normalization}
   r_i^{\{i\}}=0,\qquad c_i<0,\qquad i\in I,
\end{equation}
which is the convention in the motivating problem. Except where explicitly
indicated, the quantitative finite-grid theorem does not require
\eqref{eq:normalization}.

At stage $t\in\NN$, each player independently chooses $C_i$ or $Q_i$;
let $\theta\in\NN\cup\{\infty\}$ be the first stage with at least one
quitter, with $\theta=\infty$ if there is none. A behavior strategy for
player $i$ is the sequence $x_i=(x_i^t)_{t\ge0}$, where $x_i^t$ is the
conditional probability of quitting at $t$, provided the game has not
ended. The undiscounted (terminal) payoff is
\begin{equation}\label{eq:gamma}
 \gamma_i(x)=\EE_x\!\left[
 \one_{\{\theta<\infty\}}r_i^{S_\theta}
 +\one_{\{\theta=\infty\}}c_i\right].
\end{equation}
A profile is an undiscounted $\eps$-Nash equilibrium if for every player
$i$ and every behavior deviation $\widetilde x_i$,
\begin{equation}\label{eq:undisc-eq}
 \gamma_i(\widetilde x_i,x_{-i})-\gamma_i(x)\le\eps.
\end{equation}
The equilibrium is $\eps$-subgame-perfect if the analogous inequality
holds in the continuation game starting at every stage $t$.

For stochastic-game uniformity, give the nonabsorbing state the stage
payoff $c$, and from the first quitting stage onward give the constant
absorbing payoff $r^{S_\theta}$. If $g_i^t$ is the resulting stage-payoff
process, define
\begin{equation}\label{eq:evaluations}
 \gamma_i^T(x)=\EE_x\left[\frac1T\sum_{t=0}^{T-1}g_i^t\right],\qquad
 \gamma_i^\lambda(x)=\EE_x\left[\sum_{t\ge0}\lambda(1-\lambda)^t g_i^t\right].
\end{equation}
A \emph{uniform $\eps$-equilibrium} is one profile $x$ for which,
for all sufficiently large $T$ and all sufficiently small $\lambda>0$,
no unilateral behavior deviation improves either payoff in
\eqref{eq:evaluations} by more than $\eps$.

We use the payoff coordinate range
\begin{equation}\label{eq:D}
 D(r):=\max_{i\in I}\bigl(\max_{S\subseteq I}r_i^S-\min_{S\subseteq I}r_i^S\bigr).
\end{equation}
All payoff comparisons below are bounded by $D(r)$, regardless of the
absolute size of the payoffs.

\section{Stopping-time normal form and the tail discontinuity}\label{sec:stopping}
Set $K=\NN\cup\{\infty\}$ with the order topology of the one-point
compactification of $\NN$. A player's behavior strategy is equivalent to
a distribution $\mu_i$ of a \emph{planned quitting time} $T_i\in K$:
\begin{equation}\label{eq:hazards}
 \mu_i(t)=x_i^t\prod_{s<t}(1-x_i^s),\qquad
 \mu_i(\infty)=\prod_{s\ge0}(1-x_i^s).
\end{equation}
Conversely, as long as the denominator is positive,
\begin{equation}\label{eq:reverse-hazard}
 x_i^t=\frac{\mu_i(t)}{\mu_i(\{t,t+1,\ldots\}\cup\{\infty\})};
\end{equation}
subsequent hazards may be assigned arbitrarily when the denominator is
zero. Planned times $T_1,\ldots,T_n$ are mutually independent. Their
minimum identifies the absorption time, and the set of minimizers
identifies the quitting coalition. This construction is equivalent to
behavior strategies, including for unilateral deviations.

Let $U_i(t;\mu_{-i})$ be the expected payoff to player $i$ from the pure
planned quitting time $t\in K$, with opponents following $\mu_{-i}$.
Expected payoff is affine in a player's own distribution, so
\begin{equation}\label{eq:regret}
 \Reg_r(\mu):=\max_{i\in I}\left\{
 \sup_{t\in K}U_i(t;\mu_{-i})-\gamma_i(\mu)\right\}
\end{equation}
is precisely the maximum possible gain from \emph{any} unilateral
behavior deviation. In particular,
\begin{equation}\label{eq:g}
 g(r):=\inf_{\mu\in\prod_i\Delta(K)}\Reg_r(\mu)
\end{equation}
vanishes exactly when the game has undiscounted $\eps$-equilibria for
all $\eps>0$.

\begin{proposition}[Exact tail discontinuity]\label{prop:tail}
For arbitrary independent opponents' quitting times, let
$q_j=\mu_j(\infty)$. Under \eqref{eq:normalization},
\begin{equation}\label{eq:tail}
 \lim_{t\to\infty}\bigl(U_i(t;\mu_{-i})-U_i(\infty;\mu_{-i})\bigr)
 =(-c_i)\prod_{j\ne i}q_j.
\end{equation}
Without the normalization, the right-hand side is
$(r_i^{\{i\}}-c_i)\prod_{j\ne i}q_j$.
\end{proposition}
\begin{proof}
Fix a realization of the opponents' planned times. If at least one
opponent quits at a finite time, sufficiently late finite $t$ and
$t=\infty$ give the same payoff. If every opponent chooses $\infty$,
a finite $t$ gives the sole-quitter payoff $r_i^{\{i\}}$, whereas
$\infty$ gives $c_i$. Take expectations and apply bounded convergence;
the probability of the latter event is $\prod_{j\ne i}q_j$.
\end{proof}

The payoff kernel is continuous at pure time profiles with at least one
finite coordinate, but may be discontinuous at
$(\infty,\ldots,\infty)$. Formula \eqref{eq:tail} identifies the lost
payoff when finite quitting dates move arbitrarily far into the future.
Ordinary weak compactness of quitting-time distributions consequently
does not preserve best-response inequalities automatically.

\section{The exact-terminal-truncation attempt}\label{sec:trunc}
For $h\ge1$, let $H_h$ be the finite normal-form timing game in which
all players choose a planned time from
\begin{equation}\label{eq:Kh}
 K_h:=\{0,1,\ldots,h-1,\infty\},
\end{equation}
using the original payoff rule (in particular, the all-$\infty$ outcome
pays $c$). Such a game has a mixed Nash equilibrium by the finite-game
existence theorem. Every such mixed profile can also be embedded as a
behavior profile of the infinite quitting game.

\begin{lemma}[Boundary estimate for an exact truncation equilibrium]\label{lem:boundary}
Assume \eqref{eq:normalization}. If $\mu\in\NE(H_h)$ and
$q_i=\mu_i(\infty)$, then in the original infinite quitting game
\begin{equation}\label{eq:eta}
 \Reg_r(\mu)\le\eta(\mu),\qquad
 \eta(\mu):=\max_i(-c_i)\prod_{j\ne i}q_j.
\end{equation}
If $\infty$ is in the support of player $i$'s equilibrium distribution,
then the maximal deviation gain \emph{for that player} is exactly
$(-c_i)\prod_{j\ne i}q_j$.
\end{lemma}
\begin{proof}
Every pure deviation from $K_h$ is available in $H_h$ and hence
cannot improve player $i$'s payoff. For any finite $t\ge h$, since
all opponents' finite planned dates are less than $h$,
\begin{equation}\label{eq:tailfinite}
 U_i(t;\mu_{-i})-U_i(\infty;\mu_{-i})
 =(-c_i)\prod_{j\ne i}q_j.
\end{equation}
Since $\infty$ is an available deviation in $H_h$, the equilibrium
payoff is at least $U_i(\infty;\mu_{-i})$, which proves the upper
bound. When $\infty$ has positive equilibrium probability for player $i$,
that action is a best response in $H_h$ and has payoff exactly equal
to her equilibrium payoff; the late deviation realizes the bound.
\end{proof}

The lemma suggests choosing terminal-game equilibria whose residual
never-quit probabilities satisfy
\begin{equation}\label{eq:badcriterion}
 \inf_{h\ge1}\min_{\mu\in\NE(H_h)}
 \max_i(-c_i)\prod_{j\ne i}\mu_j(\infty)=0.
\end{equation}
If this condition held, the lemma would yield arbitrarily accurate
undiscounted equilibria. It is a \emph{sufficient} criterion, but it is
\emph{not necessary}. Section~\ref{sec:counterexample} gives a complete
three-player refutation of its universal validity.

\section{A cyclic game defeating exact terminal truncation}\label{sec:counterexample}
Consider the three-player game $\Gamma_*$ with quitting-coalition
payoffs in Table~\ref{tab:cyclic}. This is the normalized classical
cyclic-game example associated with Flesch, Thuijsman, and Vrieze
\cite{Flesch1997,AKRS2026}; its use here is as a counterexample to a
\emph{proof method}, not as a counterexample to equilibrium existence.
\begin{table}[htbp]
\centering
\caption{The cyclic three-player quitting game $\Gamma_*$.}
\label{tab:cyclic}
\begin{tabular}{@{}cl@{}}\toprule
Quitting set $S$ & Payoff vector $r^S$\\\midrule
$\varnothing$ & $(-1,-1,-1)$\\
$\{1\}$ & $(0,2,-1)$\\
$\{2\}$ & $(-1,0,2)$\\
$\{3\}$ & $(2,-1,0)$\\
$\{1,2\}$ & $(0,-1,0)$\\
$\{1,3\}$ & $(-1,0,0)$\\
$\{2,3\}$ & $(0,0,-1)$\\
$\{1,2,3\}$ & $(-1,-1,-1)$\\\bottomrule
\end{tabular}
\end{table}

\subsection{An exact cyclic subgame-perfect equilibrium}
Consider the infinite three-stage periodic behavior profile
\begin{equation}\label{eq:cycle}
 (1/2,0,0),\qquad(0,1/2,0),\qquad(0,0,1/2),
 \quad\text{repeated forever.}
\end{equation}
Here the triple at each stage lists quitting probabilities for
players $1,2,3$. The associated continuation payoffs at the beginning
of the three phases are
\begin{equation}\label{eq:cyclevalues}
 v^{(1)}=(0,1,0),\quad v^{(2)}=(0,0,1),\quad v^{(3)}=(1,0,0).
\end{equation}
Indeed,
\[
 v^{(1)}=\tfrac12 r^{\{1\}}+\tfrac12 v^{(2)}=(0,1,0),
\]
with the other two recursions following cyclically. At phase $1$,
player $1$ is indifferent between quitting and continuing (each yields
$0$); player $2$ earns $1$ by continuing and $-1/2$ by quitting;
player $3$ earns $0$ by either action. The other phases follow by
rotation of player indices.

These one-stage incentive checks are sufficient for arbitrary
deviations: under \emph{any} deviation by one player, the other two
players independently quit with probability $1/2$ at their respective
positions once per block. Their joint probability of not absorbing
the game for $k$ full blocks is $4^{-k}$. Thus any player's payoff under
any deviation differs negligibly, as $k\to\infty$, from one determined
by the first $k$ blocks. Finite-horizon backward optimality and the
one-stage checks establish all best-response inequalities.
We have proved:
\begin{proposition}\label{prop:cycle}
The infinite game $\Gamma_*$ possesses a subgame-perfect
undiscounted exact $0$-equilibrium, given by \eqref{eq:cycle}.
\end{proposition}

\subsection{A one-stage uniqueness calculation}
We now calculate \emph{all} Nash equilibria of the terminal games
$H_h$ in this example. Consider a one-stage game whose payoff, if all
continue, is the common number $v\in[-1,0)$ for each player. Write
$x_i$ for player $i$'s probability of quitting. The advantage of
quitting instead of continuing for player $1$ is
\begin{equation}\label{eq:D1}
 D_1=x_2-3x_3+x_2x_3-v(1-x_2)(1-x_3),
\end{equation}
with the analogous cyclic expressions for the other players.

\begin{lemma}[Unique one-stage equilibrium]\label{lem:onestage}
For $-1\le v<0$, the one-stage game has a unique Nash equilibrium,
which is symmetric and completely mixed, with
\begin{equation}\label{eq:onestageprob}
 x_1=x_2=x_3=p(v):=1-\frac1{\sqrt{1-v}}.
\end{equation}
Its equilibrium payoff vector is $(-p(v),-p(v),-p(v))$.
Furthermore, for \emph{any} continuation payoff vector $v\in\R^3$
(not necessarily symmetric), no Nash equilibrium of the one-stage
game has any player quit with probability $1$.
\end{lemma}
\begin{proof}
First prove the final assertion. If $x_1=1$, then player $2$'s
advantage from quitting is $D_2=2x_3-3<0$, so $x_2=0$.
Given $x_1=1,x_2=0$, player $3$'s quitting advantage is $1>0$, so
$x_3=1$. But then player $1$'s quitting advantage is $-3<0$,
a contradiction. Cyclic relabeling excludes $x_2=1$ and $x_3=1$.

Now assume $v_1=v_2=v_3=v\in[-1,0)$. If $x_1=0$, player $2$'s
quitting advantage is $x_3-v(1-x_3)>0$, forcing $x_2=1$, which is
impossible. By symmetry every $x_i$ belongs to $(0,1)$.
Set the quitting odds $z_i=x_i/(1-x_i)>0$. Dividing the
indifference equations $D_i=0$ by the positive products of the other
players' continuation probabilities yields
\begin{align}\label{eq:odds}
 v&=z_2-3z_3-z_2z_3,\\
 v&=z_3-3z_1-z_3z_1,\nonumber\\
 v&=z_1-3z_2-z_1z_2.\nonumber
\end{align}
In particular, if some $z_1\ge1$ then the second equation gives
$v=z_3(1-z_1)-3z_1\le-3$, contradicting $v\ge-1$.
Thus $0<z_i<1$ for all $i$. Put
\[
 f_v(z)=\frac{v+3z}{1-z},\qquad f_v'(z)=\frac{3+v}{(1-z)^2}>0.
\]
Equations \eqref{eq:odds} assert that
$z_2=f_v(z_3)$, $z_3=f_v(z_1)$, and $z_1=f_v(z_2)$.
A strictly increasing real map has no nontrivial finite cycle, so
$z_1=z_2=z_3$, hence $x_1=x_2=x_3=x$.
Equation \eqref{eq:D1} now reduces to
\[
 v=-\frac{x(2-x)}{(1-x)^2}=1-\frac1{(1-x)^2},
\]
which uniquely yields \eqref{eq:onestageprob}. Direct evaluation of
the pure-quitting payoff gives $-x$; by indifference, this is the
one-stage equilibrium payoff. The fully mixed profile satisfying these
equalities is indeed a Nash equilibrium.
\end{proof}

\subsection{All exact terminal equilibria and a positive regret gap}
Define recursively
\begin{equation}\label{eq:p-rec}
 p_0=1,\qquad
 p_{k+1}=1-\frac1{\sqrt{1+p_k}},\qquad k\ge0,
\end{equation}
and put
\begin{equation}\label{eq:qrec}
 q_h:=\prod_{k=1}^{h}(1-p_k).
\end{equation}

\begin{theorem}[Failure of the universal truncation criterion]\label{thm:terminal-obstruction}
For every integer $h\ge1$, the finite terminal timing game $H_h$
associated with $\Gamma_*$ has \emph{exactly one} Nash equilibrium
$\mu^{(h)}$. It has all three players quit independently with
probabilities $p_h,p_{h-1},\ldots,p_1$ at the successive stages and
has per-player payoff $-p_h$. When interpreted in the infinite quitting
game,
\begin{equation}\label{eq:qgap}
 \Reg_{\Gamma_*}(\mu^{(h)})=q_h^2>\frac14.
\end{equation}
In particular, criterion \eqref{eq:badcriterion} is false even for a
game possessing an exact subgame-perfect equilibrium.
\end{theorem}
\begin{proof}
The continuation payoffs after the last available finite date are
$(-1,-1,-1)$. By Lemma~\ref{lem:onestage}, the one-stage game at
the last date has the unique symmetric mixed equilibrium with quitting
probability $p_1$ and continuation payoff $-p_1$.

We must justify applying backward induction to \emph{all} Nash
equilibria of $H_h$, rather than just subgame-perfect equilibria.
At any continuation history reached with positive probability, a Nash
equilibrium induces a one-stage Nash equilibrium when its actual
continuation payoffs are inserted at the all-continue action profile.
By the final assertion of Lemma~\ref{lem:onestage}, no player quits
with probability $1$ at such a history. Hence the next all-continue
history also has positive probability. Induction from the initial stage
shows that \emph{every} continuation history is reached with positive
probability. Consequently, Nash equilibrium behavior is optimal at
every stage, and the one-stage uniqueness lemma can be applied
backwards. Its recursive formula is exactly \eqref{eq:p-rec}; this
proves existence and uniqueness of the specified finite equilibrium.

Each player's probability of choosing $\infty$ is $q_h>0$.
Therefore the terminal action $\infty$ lies in every player's
equilibrium support. Lemma~\ref{lem:boundary}, including its equality
case, now gives
$\Reg_{\Gamma_*}(\mu^{(h)})=q_h^2$.

It remains to prove the uniform lower bound. From \eqref{eq:p-rec},
\[
 1+p_{k-1}=\frac1{(1-p_k)^2},\qquad
 (1-p_k)(1+p_{k-1})=\frac1{1-p_k}\ge1+p_k.
\]
Multiplying and telescoping gives
\[
 q_h\ge\prod_{k=1}^h\frac{1+p_k}{1+p_{k-1}}
 =\frac{1+p_h}{2}>\frac12,
\]
because $p_h>0$. Squaring proves \eqref{eq:qgap}.
\end{proof}

\begin{remark}[What was actually disproved]
The theorem is \emph{not} a counterexample to approximate-equilibrium
existence; Proposition~\ref{prop:cycle} gives an exact equilibrium for
the same game. Nor does it show failure of arbitrary finite-prefix
approximations. It disproves the particular program of taking
\emph{exact Nash equilibria} of terminal games with an artificial
all-continue terminal payoff and expecting their infinite-game regret
to tend to zero.
\end{remark}

\subsection{Finite cyclic prefixes do approximate equilibria}
Repeat the three-stage block \eqref{eq:cycle} exactly $k$ times,
then instruct every player to continue forever. Call this profile
$x^{[k]}$. One block survives with probability $1/8$ and has
unconditional absorbing-payoff contribution $(0,7/8,0)$. Thus
\begin{equation}\label{eq:prefix-payoff}
 \gamma(x^{[k]})=
 \bigl(-8^{-k},\,1-2\cdot8^{-k},\,-8^{-k}\bigr).
\end{equation}
For an unrestricted deviator after the finite prefix, the appropriate
terminal best-response value is
$\max\{r_i^{\{i\}},c_i\}=0$, \emph{not} $-1$: such a deviator may
quit alone after everyone else has been instructed to continue.
For clarity, start the best-response backward induction at
$W=(0,0,0)$ after the last prescribed block. At a stage where only
player $j$ may quit, with probability $1/2$, the deviator $i=j$
chooses between $0$ and $W_j$; a deviator $i\ne j$ chooses between
$\tfrac12 r_i^{\{i,j\}}$ and
$\tfrac12 r_i^{\{j\}}+\tfrac12 W_i$.
Applying these recursions backwards through the stages with active
players $3,2,1$ gives
\[
 (0,0,0)\xrightarrow{j=3}(1,0,0)
 \xrightarrow{j=2}(0,0,1)
 \xrightarrow{j=1}(0,1,0).
\]
Starting instead from $(0,1,0)$ gives exactly the same three
intermediate vectors and ends again at $(0,1,0)$; thus appending
whole blocks does not change the best-response payoff vector.
Consequently,
\begin{equation}\label{eq:prefix-reg}
 \Reg_{\Gamma_*}(x^{[k]})=2\cdot8^{-k}\longrightarrow0.
\end{equation}
This demonstrates the exact distinction: arbitrary finite-support
profiles can approximate an infinite equilibrium even when exact
terminal Nash equilibria cannot.

\section{The finite-grid replacement}\label{sec:finite-grid}
The obstruction motivates optimizing regret directly over arbitrary
finite-support profiles, rather than first imposing the Nash equations
of a truncated terminal game. The next result is a quantitative
approximation of the \emph{global} infimum $g(r)$.

Fix $k\ge1$ and set $H=nk$. Let $\mathcal G_{n,k}$ denote the set of
all product distributions $\mu=\bigotimes_i\mu_i$ in which each marginal
$\mu_i$ is the empirical distribution of $k$ (not necessarily distinct)
points in
\begin{equation}\label{eq:grid-space}
 K_H=\{0,1,\ldots,H-1,\infty\}.
\end{equation}
These are exactly the marginals whose probability masses are integer
multiples of $1/k$. There are
\begin{equation}\label{eq:gridcount}
 |\mathcal G_{n,k}|=\binom{(n+1)k}{k}^{n}
\end{equation}
independent profiles. Define
\begin{equation}\label{eq:mk}
 m_k(r)=\min_{\mu\in\mathcal G_{n,k}}\Reg_r(\mu).
\end{equation}
We stress that \eqref{eq:mk} uses unrestricted infinite-game regret.

\begin{theorem}[Quantitative finite-grid approximation]\label{thm:grid}
For every $n\ge2$, every finite-player quitting-game payoff table
$r$, and every positive integer $k$,
\begin{equation}\label{eq:grid-main}
 0\le m_k(r)-g(r)\le\frac{(2n-1)D(r)}k.
\end{equation}
Equivalently,
\begin{equation}\label{eq:grid-interval}
 \max\!\left\{0,m_k(r)-\frac{(2n-1)D(r)}k\right\}
 \le g(r)\le m_k(r).
\end{equation}
No assumption such as \eqref{eq:normalization} is needed.
\end{theorem}
\begin{proof}
\emph{Step 1: midpoint-quantile rounding.}
Fix any profile $\mu=\bigotimes_i\mu_i$. For each $i$ and
$\ell\in\{1,\ldots,k\}$, choose the leftmost $t_{i,\ell}\in K$
whose cumulative distribution function reaches
$u_\ell=(2\ell-1)/(2k)$, and define
\[
 \nu_i=\frac1k\sum_{\ell=1}^k\delta_{t_{i,\ell}}.
\]
A quantile may equal $\infty$. Standard midpoint rounding gives,
for every finite $t$, as well as for every left limit,
\begin{equation}\label{eq:cdf-control}
 \bigl|F_{\nu_i}(t)-F_{\mu_i}(t)\bigr|\le\frac1{2k}.
\end{equation}
The inequality also applies to the total finite mass, by taking
$t\to\infty$.

\emph{Step 2: the payoff is a three-level function of one stopping
time.} Fix every stopping time except that of player $j$ and inspect
any coordinate $i$. If the earliest among the fixed times is a finite
$a$, and $J$ is the set of fixed players quitting at $a$, then as a
function of $s=T_j$ the payoff is
\[
 f(s)=\begin{cases}
 A=r_i^{\{j\}},&s<a,\\
 B=r_i^{J\cup\{j\}},&s=a,\\
 C=r_i^J,&s>a.
 \end{cases}
\]
Let $\Delta F=F_{\nu_j}-F_{\mu_j}$. The change in the expectation
of $f$ is exactly
\[
 (A-B)\Delta F(a-)+(B-C)\Delta F(a).
\]
By \eqref{eq:D} and \eqref{eq:cdf-control}, its absolute value is
at most $D(r)/k$. If all fixed times equal $\infty$, there are only
two payoff levels, and the same bound holds. Averaging over the fixed
stopping times preserves the bound.

\emph{Step 3: all unilateral deviations are controlled.}
Replace the $n$ marginals successively. The bound just established
gives, simultaneously for all players,
\begin{equation}\label{eq:gamma-round}
 |\gamma_i(\nu)-\gamma_i(\mu)|\le\frac{nD(r)}k.
\end{equation}
For any fixed pure deviation $t\in K$, only the other $n-1$
marginals change, and the same estimate is uniform in $t$:
\begin{equation}\label{eq:U-round}
 \sup_{t\in K}\bigl|U_i(t;\nu_{-i})-U_i(t;\mu_{-i})\bigr|
 \le\frac{(n-1)D(r)}k.
\end{equation}
Combining \eqref{eq:gamma-round} and \eqref{eq:U-round} yields
\begin{equation}\label{eq:reg-round}
 \Reg_r(\nu)\le\Reg_r(\mu)+\frac{(2n-1)D(r)}k.
\end{equation}

\emph{Step 4: date compression.}
The union of the finite supports of the $\nu_i$ has at most $nk$
points. Let the distinct finite support dates be
$s_0<\cdots<s_{L-1}$, with $L\le nk$. Replace the date $s_a$ by $a$
for every player, leaving $\infty$ unchanged, and call the product
profile $\widehat\nu$. Every realized order and every tie of planned
times is preserved, so
\begin{equation}\label{eq:compress-gamma}
 \gamma_i(\widehat\nu)=\gamma_i(\nu).
\end{equation}
For deviations, each new finite quitting time $a<L$ has the same
payoff as the original deviation $s_a$; every new finite time
$t\ge L$ has the same payoff as the original deviation
$s_{L-1}+1$; and the deviation $\infty$ is unchanged. If $L=0$,
the same conclusion follows by comparing immediate finite quitting
with never quitting. Thus compression may delete old deviation
opportunities but cannot create a new payoff level, and
\begin{equation}\label{eq:compress-reg}
 \Reg_r(\widehat\nu)\le\Reg_r(\nu).
\end{equation}
The profile $\widehat\nu$ belongs to $\mathcal G_{n,k}$.
Consequently, for arbitrary $\mu$,
\[
 m_k(r)\le\Reg_r(\widehat\nu)
 \le\Reg_r(\mu)+\frac{(2n-1)D(r)}k.
\]
Take the infimum over $\mu$ to obtain the upper bound in
\eqref{eq:grid-main}. The lower bound follows from
$\mathcal G_{n,k}\subseteq\prod_i\Delta(K)$.
\end{proof}

\begin{remark}[Neither correlation nor a compactness shortcut]
Each rounded marginal remains independent of the others. Date
compression is the same deterministic relabeling for every player;
it preserves independence. No public signal is introduced. The
argument avoids a hidden appeal to continuity at the all-$\infty$
profile by bounding payoff errors before compressing dates.
\end{remark}

\section{Exact grid regret and finite nonexistence certificates}\label{sec:certificates}
For a profile supported on $K_H$, the deviation payoff is constant
for finite quitting times $t\ge H$. In full generality, the set of
candidate deviations reduces to
\begin{equation}\label{eq:finite-deviation-list}
 \{0,1,\ldots,H,\infty\}.
\end{equation}
Under \eqref{eq:normalization}, set $q_j=\mu_j(\infty)$.
The exact late-deviation identity is
\begin{equation}\label{eq:late-general}
 U_i(H;\mu_{-i})=U_i(\infty;\mu_{-i})
 -c_i\prod_{j\ne i}q_j.
\end{equation}
Since $-c_i>0$, the late finite deviation weakly dominates never
quitting, and therefore
\begin{equation}\label{eq:finite-regret-formula}
 \Reg_r(\mu)=\max_i\Bigg[
 \max\left\{U_i(0;\mu_{-i}),\ldots,
 U_i(H-1;\mu_{-i}),
 U_i(\infty;\mu_{-i})-c_i\prod_{j\ne i}q_j\right\}
 -\gamma_i(\mu)\Bigg].
\end{equation}
Importantly, the deviation immediately \emph{after} the last
date cannot be omitted.

\begin{corollary}[Finite certificate of failure of approximate equilibrium]\label{cor:certificate}
Set $C(r)=(2n-1)D(r)$. The following are equivalent:
\begin{align}\label{eq:finite-iff}
 &g(r)>0;\\
 &\exists\,k\ge1\text{ such that }m_k(r)>\frac{C(r)}k.\notag
\end{align}
If an exact calculation establishes $m_k(r)\ge L>C(r)/k$, then
\begin{equation}\label{eq:cert-low}
 \Reg_r(\mu)\ge L-\frac{C(r)}k>0
 \quad\text{for \emph{every} infinite strategy profile }\mu.
\end{equation}
\end{corollary}
\begin{proof}
Theorem~\ref{thm:grid} proves the implication from the second
statement to the first and \eqref{eq:cert-low}. Conversely, if
$g(r)>0$, choose $k$ with $C(r)/k<g(r)$; since
$m_k(r)\ge g(r)$, the second statement follows.
\end{proof}

\begin{proposition}[Rational counterexamples would suffice]\label{prop:lipschitz}
For two games $r,\widetilde r$ with the same player set,
\begin{equation}\label{eq:Lipschitz}
 |g(r)-g(\widetilde r)|\le2\|r-\widetilde r\|_\infty.
\end{equation}
Thus if a counterexample exists with real payoffs, there is also a
counterexample with rational payoffs. Rational approximation may be
chosen to preserve the exact singleton-payoff normalization and the
strict signs $c_i<0$.
\end{proposition}
\begin{proof}
For fixed $\mu$, replacing $r$ by $\widetilde r$ changes every
actual payoff and every pure-deviation payoff by at most
$\|r-\widetilde r\|_\infty$. Therefore it changes
$\Reg_r(\mu)$ by at most twice this amount. Infima over $\mu$
preserve the same Lipschitz bound. Choose a rational table sufficiently
close to a hypothetical $r$ with $g(r)>0$.
\end{proof}

\begin{corollary}[Equivalent universal finite inequalities]\label{cor:positive}
Every $n$-player quitting game has undiscounted $\eps$-equilibria
for every $\eps>0$ if and only if
\begin{equation}\label{eq:unproved}
 m_k(r)\le\frac{(2n-1)D(r)}k
 \qquad\text{for every payoff table $r$ and every }k\ge1.
\end{equation}
\end{corollary}
\begin{proof}
If all such games have $g(r)=0$, Theorem~\ref{thm:grid} gives
\eqref{eq:unproved}. Conversely, \eqref{eq:unproved} and
$0\le g(r)\le m_k(r)$ imply $g(r)=0$ upon letting $k\to\infty$.
\end{proof}

The inequalities \eqref{eq:unproved} are \emph{not} proved here.
Theorem~\ref{thm:grid} gives only
$m_k(r)\le g(r)+C(r)/k$. Replacing the right-hand side with
$C(r)/k$ without another argument would assume the conclusion of
the open conjecture. At fixed $n,k$, each grid profile's regret is the
maximum of finitely many linear expressions in the payoff table,
with rational coefficients. A complete lower-bound certificate for a
fixed rational game is therefore finite and exactly checkable. It
may, however, require an astronomically large enumeration.

\section{From finite undiscounted profiles to uniform approximations}\label{sec:uniform}
The reduction has a direct uniform-equilibrium interpretation.
Let $\Reg_T(\mu)$ and $\Reg_\lambda(\mu)$ be the maximum unilateral
gains under the $T$-stage average and $\lambda$-discounted evaluations
in \eqref{eq:evaluations}, respectively.

\begin{theorem}[Uniform error bounds for finite-support profiles]\label{thm:uniform}
If each player's stopping-time distribution is supported on
$K_H=\{0,\ldots,H-1,\infty\}$, then
\begin{equation}\label{eq:uniform-bound}
 \Reg_T(\mu)\le\Reg_r(\mu)+\frac{2D(r)H}{T},
 \qquad
 \Reg_\lambda(\mu)\le\Reg_r(\mu)+2D(r)H\lambda.
\end{equation}
\end{theorem}
\begin{proof}
Under the prescribed profile, either absorption occurs at a stage
less than $H$, or nobody ever quits. Therefore its $T$-average payoff
differs from the terminal payoff by at most $D(r)H/T$, and the
discounted payoff differs by at most
$D(r)[1-(1-\lambda)^H]\le D(r)H\lambda$.

The same bounds apply to any pure deviation to a quitting time
$t\le H$ or $t=\infty$: absorption then either occurs before $H$
or never occurs. If a deviator instead chooses finite $t>H$, all
opponents who quit do so before $H$ and thereby make the late
deviation irrelevant. Conditional on all opponents choosing $\infty$,
the evaluated payoff is a convex combination of $c_i$ and
$r_i^{\{i\}}$, with the weight on $r_i^{\{i\}}$ no larger than it is
at time $H$. For either evaluation, such a late deviation has a
payoff no greater than the maximum of quitting at $H$ and never
quitting. Hence the evaluated best-response payoff exceeds the
undiscounted best-response payoff by at most the respective one-sided
bound. Subtract the lower bound for the prescribed profile's own
evaluated payoff. Mixed deviations are convex combinations of pure
deviations, completing the proof.
\end{proof}

\begin{corollary}[Ordinary and uniform approximate existence are equivalent]\label{cor:equivalence}
For any finite-player quitting game in the sense of
Section~\ref{sec:model}, the following are equivalent:
\begin{enumerate}[label=(\roman*)]
\item for every $\eps>0$ there is an undiscounted $\eps$-equilibrium;
\item for every $\eps>0$ there is a uniform $\eps$-equilibrium.
\end{enumerate}
\end{corollary}
\begin{proof}
(i)$\Rightarrow$(ii): fix $\eps>0$ and an undiscounted profile of
regret below $\eps/3$. By Theorem~\ref{thm:grid}, or directly by
quantile rounding and date compression, it can be approximated by
a finite-support profile whose regret is less than $2\eps/3$.
Theorem~\ref{thm:uniform} makes this profile a uniform $\eps$-
equilibrium for sufficiently large $T$ and small $\lambda$.

(ii)$\Rightarrow$(i): for any fixed behavior profile and any fixed
unilateral deviation, bounded convergence gives convergence of the
long-run averages to the undiscounted payoff. A uniform $\eps/2$-
equilibrium satisfies the average-payoff deviation inequality for all
large $T$. Passing to the limit separately for each deviation yields
an undiscounted $\eps/2$-equilibrium, hence an $\eps$-equilibrium.
\end{proof}

No corresponding assertion about subgame perfection follows just
from the quantile approximation: rounding and compression do not
preserve incentive constraints at every continuation history.

\section{Exact finite-grid calculations for the cyclic example}\label{sec:computation}
The finite-grid theorem is constructive in a theoretical sense.
For a rational payoff table, \eqref{eq:finite-regret-formula} is a
rational expression on every grid profile, so $m_k(r)$ is an exactly
computable rational minimum by exhaustive enumeration. A sampled or
local numerical minimum is \emph{not} a certificate of nonexistence.

For the payoff table in Table~\ref{tab:cyclic}, exact rational
enumeration gives the following small-grid results:
\begin{center}
\begin{tabular}{@{}cccc@{}}\toprule
$k$ & $H=nk$ & Number of grid profiles & Exact $m_k(\Gamma_*)$\\\midrule
$1$ & $3$ & $64$ & $1$\\
$2$ & $6$ & $21{,}952$ & $1/8$\\\bottomrule
\end{tabular}
\end{center}
One minimizing profile for $k=2$ is the independent product of
\begin{equation}\label{eq:grid-min-example}
 T_1\sim\tfrac12\delta_0+\tfrac12\delta_3,\quad
 T_2\sim\tfrac12\delta_1+\tfrac12\delta_4,\quad
 T_3\sim\tfrac12\delta_2+\tfrac12\delta_4.
\end{equation}
Direct enumeration of the $2^3$ joint planned-time outcomes and of all
best-response date classes yields
\begin{equation}\label{eq:grid-min-payoff}
 \gamma(\mu)=(0,9/8,-1/8),\qquad
 \bigl(\sup_t U_i(t;\mu_{-i})\bigr)_{i=1}^3=(0,5/4,0),
\end{equation}
so the unrestricted regret is $1/8$. Since
$D(\Gamma_*)=3$ and $(2n-1)D=15$, the theorem's bound
$m_2-15/2$ gives no positive lower bound on $g(\Gamma_*)$.
In fact Proposition~\ref{prop:cycle} shows that
$g(\Gamma_*)=0$. The computations are a check of the explicit
finite formulation, not evidence of nonexistence.

The exhaustive enumeration has a simple implementation-independent
specification. Enumerate, for every player, the
$\binom{H+k}{k}$ multisets of $k$ dates from $K_H$. For each product
profile and every $i$, compute
\begin{align*}
 \gamma_i(\mu)
 &=\sum_{(t_1,\ldots,t_n)\in K_H^n}
   \left(\prod_j\mu_j(t_j)\right) r_i^{S(t_1,\ldots,t_n)},\\
 U_i(t;\mu_{-i})
 &=\sum_{t_{-i}\in K_H^{n-1}}
   \left(\prod_{j\ne i}\mu_j(t_j)\right)r_i^{S(t,t_{-i})},
\end{align*}
where $S(\cdot)$ is the set of indices attaining the least finite
date, and $r^\varnothing=c$ when every date is infinite. Evaluate
the $H+2$ candidate deviations in \eqref{eq:finite-deviation-list}
and minimize \eqref{eq:regret}. This is a finite exact certificate
procedure involving no numerical tolerances.

\section{Absorption paths, the residual obstacle, and conclusions}\label{sec:discussion}
The absorption-path approach of Ashkenazi-Golan, Krasikov, Rainer,
and Solan \cite{AKRS2024} compactifies absorbing behavior profiles
using accumulated absorption probability in place of physical time.
Its sequential-perfectness conditions retain both jump behavior
(positive simultaneous quitting probabilities at discrete events)
and continuous components (vanishing quitting rates). The connection
between such paths and limits of approximate equilibria depends on
the exclusions of the elementary immediate-absorption and
never-absorption cases, as specified precisely in
\cite{AKRS2024,AKRS2026}. The 2026 APS treatment develops an operator
for a restricted class of paths in which one player is active at a
time \cite{AKRS2026}; it is not a proof of universal ordinary
equilibrium existence.

The main lessons of the consolidated attempt can be stated without
ambiguity:
\begin{enumerate}[label=\arabic*.]
\item Proposition~\ref{prop:tail} pinpoints the discontinuity that
makes a naive limit of finite terminal-game equilibria unsafe.
\item Lemma~\ref{lem:boundary} is a correct and useful sufficient
error estimate, but Theorem~\ref{thm:terminal-obstruction} proves that
forcing this estimate to vanish along \emph{exact} terminal-game
Nash equilibria is not a viable general argument.
\item Arbitrary finite-support profiles avoid that obstruction.
Theorem~\ref{thm:grid} approximates the globally minimal regret
with a computable, explicit $O(1/k)$ error bar; the finite certificate
Corollary~\ref{cor:certificate} is the corresponding negative-
resolution criterion.
\item Corollary~\ref{cor:positive} isolates the unproved universal
inequality \eqref{eq:unproved}. Establishing it for every payoff table,
or finding one finite verified violation, would respectively give
a positive solution or a genuine counterexample.
\end{enumerate}

\noindent\textbf{Status.} No proof of \eqref{eq:unproved}, and no
counterexample meeting Corollary~\ref{cor:certificate}, is given here.
Thus the general existence problem remains unresolved by the
arguments in this manuscript. The finite reductions are presented as
transparent mathematical research attempts; no claim is made that
they were previously unknown or that they constitute a solution.

\begin{thebibliography}{99}
\bibitem{Flesch1996}
J.~Flesch, F.~Thuijsman, and O.~J.~Vrieze (1996).
Recursive repeated games with absorbing states.
\emph{Mathematics of Operations Research} \textbf{21}, 1016--1022.

\bibitem{Flesch1997}
J.~Flesch, F.~Thuijsman, and K.~Vrieze (1997).
Cyclic Markov equilibria in stochastic games.
\emph{International Journal of Game Theory} \textbf{26}, 303--314.

\bibitem{SolanVieille2001}
E.~Solan and N.~Vieille (2001).
Quitting games.
\emph{Mathematics of Operations Research} \textbf{26}(2), 265--285.
\url{https://doi.org/10.1287/moor.26.2.265.10549}.

\bibitem{SolanSolan2020}
E.~Solan and O.~N.~Solan (2020).
Quitting games and linear complementarity problems.
\emph{Mathematics of Operations Research} \textbf{45}(2), 434--454.
\url{https://doi.org/10.1287/moor.2019.0996}.

\bibitem{AKRS2024}
G.~Ashkenazi-Golan, I.~Krasikov, C.~Rainer, and E.~Solan (2024).
Absorption paths and equilibria in quitting games.
\emph{Mathematical Programming} \textbf{203}, 735--762
(first published online 22 April 2022).
\url{https://doi.org/10.1007/s10107-022-01807-6}.

\bibitem{AKRS2026}
G.~Ashkenazi-Golan, I.~Krasikov, C.~Rainer, and E.~Solan (2026).
The APS approach for undiscounted quitting games.
\emph{International Journal of Game Theory} \textbf{55}, Article 19
(published online 2 April 2026).
\url{https://doi.org/10.1007/s00182-026-00982-6}.
\end{thebibliography}

\end{document}
